\chapter{Conventions and constructions}
\label{chap:library}

The manual treats ordinary characters, Brauer characters, reduction and
blocks, and the weights defined from ordinary characters of normaliser
quotients.

Three kinds of construction recur in these subjects: passage to a normal
subgroup or quotient, passage from representations to their characters or
Grothendieck classes, and transport under a group action. The chapters
develop these constructions together with the compatibility statements
needed to combine them.

Throughout the manual, groups are finite unless a statement says otherwise.
The modular prime is $\ell$, $k$ has characteristic $\ell$, and $K$ has
characteristic zero. A discrete valuation ring is a local principal ideal
domain that is not a field. Coefficient rings, splitting conditions,
algebraic closure and choices of roots are specified where needed. The
source files often use the letter $p$ for the modular prime. In the Conlon
argument, $p$ denotes a potentially different prime used for permutation
lattices.

\section{From a result to its source}
\label{chap:library:use}

Mathlib supplies the representation types \lean{Representation} and
\lean{FDRep}, the ordinary trace character \lean{Representation.character},
and the underlying group theory, linear algebra and category theory.
ModularRep defines the finite root correspondences, Brauer characters and
exact Grothendieck group used here. The implementation paragraphs explain
which steps use Mathlib results and identify the definitions and proofs
in the modules of this package.

The accompanying module index, linked from the package reading guide,
lists the library and formalisation modules, their declarations and stated
assumptions. Its entries are classified by role and status.
Examples in this manual illustrate the stated mathematical
results.

For local quotient arguments, one can begin with
\begin{samepage}
\begin{verbatim}
import ModularRep.PCore
import ModularRep.NormalCoreTransport
\end{verbatim}
\end{samepage}
The normaliser quotient results apply to a surjective homomorphism whose
kernel lies in the relevant subgroup. The induced isomorphism of
normaliser quotients determines the pullback of representations.
For reduction, the imports
\begin{samepage}
\begin{verbatim}
import ModularRep.BrauerDecompositionMap
import ModularRep.KZeroTwist
\end{verbatim}
\end{samepage}
provide the decomposition homomorphism for the chosen modular system and
root correspondence. Its hypotheses relate reduction of a stable lattice
to Brauer characters. The naturality theorems use the same choices. The modules
\begin{samepage}
\begin{verbatim}
import Formalisation.FibreTransport
import Formalisation.IndexedAssembly
\end{verbatim}
\end{samepage}
construct equivariant bijections of finite sets from equivariant
bijections of their fibres. The chapters state the hypotheses of each
construction.

\section{Hypotheses and proofs}

The mathematical exposition includes standard results used by the library.
An implementation paragraph says whether the corresponding result is
proved in Lean or supplied as an input. In particular,
Hypotheses~\ref{lib:blocks:catalogue} and
\ref{lib:blocks:defect-hypotheses} collect standard published facts
used as inputs to the Lean arguments. Other hypotheses express
compatibility of chosen data, as in
Hypothesis~\ref{lib:reduction:character-hypothesis}.
For example, multiplicativity of
the central Brauer map follows from cancellation of nontrivial
$\ell$-group orbits. The construction of the decomposition homomorphism
uses an explicit hypothesis relating stable reduction to Brauer
characters. The existence of a stable lattice and the deductions from
that compatibility hypothesis are proved separately.

An equivariant bijection of finite sets does not by itself preserve a
specified relation. To construct a bijection preserving a relation from
bijections of the fibres, the latter must preserve the relation and the
relation must be invariant under the acting group. Character and block
relations require their corresponding compatibility hypotheses.

Finite computations also require a mathematical interpretation. A matrix
rank calculation, for example, counts irreducible Brauer characters only
after its rows have been identified with character restrictions and the
required spanning and independence statements have been established.
Chapter~\ref{chap:library-computations} states these hypotheses and proves
the resulting rank formulas.
